% ---------------- Abstract ----------------
\chapter*{Abstract}
\addcontentsline{toc}{chapter}{Abstract}

Tobramycin is a polycationic aminoglycoside antibiotic on the WHO Model List of Essential Medicines, indispensable against \textit{Pseudomonas aeruginosa} in cystic fibrosis and nosocomial Gram-negative infections, yet available only by intravenous and inhaled routes. Although commonly labelled BCS class~IV, a structured analysis of solubility (94--100~mg/mL in water, 20--200$\times$ the dose criterion over pH~1.2--6.8) and permeability evidence (apparent $P_{\mathrm{app}}<1\times10^{-6}$~cm/s; absolute oral bioavailability $\approx$1--2\%) demonstrates that tobramycin is unambiguously \textbf{BCS class III}: the sole barrier to oral absorption is permeability. This reclassification redirects the formulation problem from solubility enhancement to permeability enhancement, for which a validated toolbox exists: hydrophobic ion pairing (HIP; logP shift from $-2.9$ to $\approx+1.6$, a 1500-fold improvement), self-emulsifying drug delivery systems (SEDDS), polymeric and colloidal nanoparticles, and gastrointestinal permeation enhancers (sodium caprate C$_{10}$, SNAC).

Building on this synthesis, a two-compartment population pharmacokinetic model (CL = 5.5~L/h, $V_1$ = 17~L, $Q$ = 2.4~L/h, $V_2$ = 16~L, $t_{1/2}$ = 2.5~h; \citet{Li2021}) was coupled to a transit-compartment absorption model and a Hill-type pharmacodynamic driver. Ten formulation design variables --- modified logP, particle size, SEDDS composition (surfactant/cosurfactant/oil), permeation-enhancer concentration, polymer loading, chitosan and enteric coatings, and dose --- were optimized by a classical genetic algorithm (population 100, 200 generations) and, subsequently, by NSGA-II (population 200, 300 generations) against four objectives: maximize $\F$, maximize $\CMIC$, maximize $\AMIC$, minimize dose. The Pareto front contained 200 non-dominated formulations. The recommended candidate (TOB-161; 550~mg enteric-coated HIP--SEDDS capsule, 50~nm polymer nanoparticles, chitosan coating, 50~mM C$_{10}$, surfactant:co-surfactant:oil = 50:30:20, polymer 30\%) achieves $\F = 34\%$ --- a 23-fold enhancement over the native molecule --- with $C_{\max} = 9.5$~mg/L ($\CMIC = 9.5$ at MIC~1~mg/L, exceeding the efficacy threshold of 8), $\AMIC = 32$, trough 0.16~mg/L, and $\fmic = 30\%$.

Uncertainty quantification (Monte Carlo, $n=200$; cross-validation with $\pm20\%$ parameter perturbation, $n=26$; one-at-a-time sensitivity analysis) identified dose and modified logP as the dominant sensitivity drivers (indices 0.42 and 0.34). A methodology audit then rebuilt the gene-to-permeability mapping as a \emph{cited, bounded, emergent} enhancement model: every component factor carries a literature range and a graded uncertainty, the logP gene is capped at the measured hydrophobic ion-pair complex, and the emergent multiplier caps at $\times$125. Two complete, standardized evaluations --- identity card, single dose, true seven-day repeated dosing, PK/PD targets, virtual population, renal impairment, food effect, component sensitivity, IV comparison --- demonstrate that both candidate formulations attain the peak and exposure targets at steady state: the re-evaluated exploratory winner (TOB-161-A, $\times$73.9, 550.6 mg once daily) lands inside the 80--120 AUC/MIC band (87.3), and the in-loop optimum (TOBP-001, $\times$126.3, 1000 mg) exceeds it with a 467 mg in-band option, both with troughs at the 1 mg/L monitoring line and mass balance verified to 0.1\%. The direct experimental test --- Caco-2/Ussing permeability of the hydrophobic ion-pair complex across ion-pair loading --- is defined as the first roadmap item. A manufacturing route of ten unit operations was costed at $\approx$\$4.30 per dose (feasibility grade B), and a 505(b)(2) regulatory pathway of approximately seven years and \$4.8~million was delineated, including preclinical toxicology of nanoparticles, mucosal tolerance of C$_{10}$, and chitosan quality control.

This thesis resolves the BCS classification controversy, delivers a reproducible, fully traceable in silico framework --- from literature parameters through evolutionary optimization to manufacturing and regulatory feasibility --- and provides ten ranked, costed, and risk-assessed candidates for the first oral tobramycin. It thereby de-risks the experimental program (Caco-2/Ussing permeability, HIP--SEDDS prototype formulation, first-in-human Phase~I) required to convert an IV-only essential antibiotic into an orally available therapeutic option.

\bigskip
\noindent\textbf{Keywords:} tobramycin; BCS class III; oral bioavailability; hydrophobic ion pairing; SEDDS; permeation enhancers; nanoparticles; PBPK; genetic algorithm; NSGA-II; Pareto optimization; cystic fibrosis.
